\documentclass[letterpaper,9pt]{article}

\usepackage{fontawesome5}
\usepackage{latexsym}
\usepackage[empty]{fullpage}
\usepackage{titlesec}
\usepackage{marvosym}
\usepackage[usenames,dvipsnames]{color}
\definecolor{teal}{RGB}{0,128,128}
\usepackage{verbatim}
\usepackage{enumitem}
\usepackage[hidelinks]{hyperref}
\usepackage{fancyhdr}
\usepackage[english]{babel}
\usepackage{tabularx}
\usepackage{enumitem} 
\input{glyphtounicode}
% \usepackage{fontspec}
\usepackage{graphicx}
\graphicspath{ {./images/} }
\usepackage{tikz}

\pagestyle{fancy}
\fancyhf{} % clear all header and footer fields
\fancyfoot{}
\renewcommand{\headrulewidth}{0pt}
\renewcommand{\footrulewidth}{0pt}

% Adjust margins
\addtolength{\oddsidemargin}{-0.5in}
\addtolength{\evensidemargin}{-0.5in}
\addtolength{\textwidth}{1in}
\addtolength{\topmargin}{-0.5in}
\addtolength{\textheight}{1.0in}

\urlstyle{same}

\raggedbottom
\raggedright
\setlength{\tabcolsep}{0in}
\setlist[itemize]{nosep, topsep=0pt, partopsep=0pt, parsep=0.2em, itemsep=0.1em}


% Formatting sections
\titleformat{\section}{
  \vspace{-4pt}\raggedright\Large
}{}{0em}{}[\color{black}\titlerule\vspace{-5pt}]

% \setmainfont[
% Extension = .ttf,
%   UprightFont = fonts/*-Regular,
%   BoldFont = fonts/*-Bold.ttf,
%   ItalicFont = fonts/*-Italic.ttf,
%   BoldItalicFont = fonts/*-BoldItalic.ttf
% ]{Calibri}

\begin{document}
\include{custom-commands}

% plaksha logo
\begin{tikzpicture}[remember picture, overlay]
    \node[anchor=north east, xshift=-0.4cm, yshift=-0.4cm] at (current page.north east) {
        \includegraphics[width=3.7cm]{logo_latest.png}
    };
\end{tikzpicture}

{\Huge \textbf{\textcolor{teal}{Avantika Bansal}}}

\vspace{0.2em}

\noindent

\begin{tabularx}{\textwidth}{@{}X@{}}
\href{mailto:avantika.bansal.ug23@plaksha.edu.in}{avantika.bansal.ug23@plaksha.edu.in} \textbar\ 
+91 9829011528 \\
\faLinkedin\ \href{https://www.linkedin.com/in/avantika-bansal-9253b3278/}{\textbf{LinkedIn}} \textbar\ 
\faGithub\ \href{https://github.com/AvantikaBnsl}{\textbf{GitHub}}
\end{tabularx}


% \vspace{0.3em}

% \section*{Education}
\noindent
\begin{center}
    \rule[0.5ex]{0.4\linewidth}{0.5pt}
    \hspace{1em}
    \textbf{EDUCATION}
    \hspace{1em}
    \rule[0.5ex]{0.4\linewidth}{0.5pt}
\end{center}

\begin{tabularx}{\textwidth}{@{} l l X r @{}}
{\small
\textbf{BTech} - Biological Systems Engineering\hspace{4em} & Plaksha University, Mohali, India\hspace{5em} & 8.93 & 2027 \\[0.5em]
\textbf{12th} - CBSE & MGD Girls’ School, Jaipur, India\hspace{4em} & 94.8\% & 2023 \\[0.5em]
\textbf{10th} - CBSE & MGD Girls’ School, Jaipur, India\hspace{4em} & 93.2\% & 2021 \\
}
\end{tabularx}

\vspace{0.2em}



% \section*{Experience}
\noindent
\begin{center}
    \rule[0.5ex]{0.34\linewidth}{0.5pt}
    \hspace{0.5em}
    \textbf{RESEARCH EXPERIENCE}
    \hspace{0.5em}
    \rule[0.5ex]{0.34\linewidth}{0.5pt}   
    \vspace{0.3em}
\end{center}

\begin{itemize}[leftmargin=0pt, parsep=0pt, itemsep=0.4em]

% \item[] \textbf{Experimental investigations on fungal microbe proliferation in agricultural settings} \textbar\ \textit{Dr. Navjot Kaur} \hfill Aug 2026 -- Present
% {\small
% \begin{itemize}[leftmargin=1em, label=\textbullet, itemsep=1pt, parsep=1pt, topsep=1pt]
% \item Designed experiments to investigate survival, proliferation, and population dynamics of beneficial and pathogenic fungi in soil–plant ecosystems.
% \item Cultured \textit{Trichoderma} strains and \textit{Alternaria alternata}, and standardized spore concentrations for controlled microbial inoculation.
% \item Established tomato soil–plant–microbe systems with controlled soil treatments and two microbial inoculum concentrations.
% \item Conducted time-point soil sampling and microbial population enumeration to quantify fungal survival and proliferation.
% \item Analysed microbial growth, plant growth, and disease parameters to identify factors influencing fungal proliferation and inform product formulation strategies.
% \end{itemize}
% }


\item[] \textbf{Research Intern, Amgen Scholars Program} \hfill Jun 2026 -- Aug 2026\\
\textbf{Computational Aptamer Optimization for Electrochemical Biosensor} \textbar\ \textit{Dr. Abhishek Srivastava, IIIT-H}
{\small
\begin{itemize}[leftmargin=1em, label=\textbullet, itemsep=0.5pt, parsep=0pt, topsep=1pt]
\item Established a \textbf{ZnO nanorod based electrochemical sensor} for rapid detection of \textbf{Tuberculosis}.
\item Fabricated and characterized \textbf{ZnO nanorod-modified IDEs using hydrothermal synthesis} for electrochemical biosensing.
\item Developed a computational pipeline for aptamer optimization using \textbf{2D and 3D modelling} with RNA structure simulation tools.
\item Extracted \textbf{sequence, structural, physicochemical, docking features and binding affinity} for aptamer–protein pairs.  
\item Presented at the\textbf{ Amgen Asia Symposium}, University of Tokyo, engaging with 75+ researchers from across Asia and Australia.
\end{itemize}
}

\vspace{0.35em}

\item[] \textbf{Assessing Fungal Microbiome to Investigate Soil Carbon Sequestration} \textbar\ \textit{Dr. Navjot Kaur} \hfill Jul 2025 -- Present
{\small
\begin{itemize}[leftmargin=1em, label=\textbullet, itemsep=0.5pt, parsep=0pt, topsep=1pt]
\item \textbf{Characterized 3 economically affecting fungal pathogens} on a molecular level for species level identification.  
\item Performed fungal culturing, microscopy and genomic DNA extraction.
\item Executed \textbf{PCR amplification and restriction digestion} for species identification with Sanger sequencing. 
\item Analyzed fungal sequences using \textbf{phylogenetic trees with MEGA} alignment tool to validate species identity and similarity.
\item Conducted a comparative analysis for \textbf{ITS and species specific primers} obtaining \>95\% alignment with reference sequences.
\end{itemize}
}

\vspace{0.35em}

\item[] \textbf{Algal Facades for Sustainable Buildings} \textbar\ \textit{Prof. Malini Balakrishnan, Prof. Prashanth Kumar} \hfill Jan 2025 -- May 2026
{\small
\begin{itemize}[leftmargin=1em, label=\textbullet, itemsep=0.5pt, parsep=0pt, topsep=1pt]
\item Developed and analyzed efficiency of \textbf{algal photobioreactors} to evaluate energy-saving potential.
\item Ran \textbf{15+ building-energy simulations} varying facade parameters with DesignBuilder software to model performance.
\item Cultured algae and carried out media preparation, inoculation, sampling, and growth-curve analysis to characterize algal growth.
\item Designed and operated a lab scale photobioreactor to \textbf{experimentally validate algal growth and nutrient absorption}.
\item Conducted \textbf{life-cycle assessment (LCA)} using OpenLCA to evaluate environmental impacts of the algal facade vs blinds.
\item Found energy savings value of upto \textbf{-20 Global Warming Potential} units for facade operation.
\end{itemize}
}

% \section*{Publications}
\noindent
\begin{center}
    \rule[0.5ex]{0.395\linewidth}{0.5pt}
    \hspace{0.5em}
    \textbf{PUBLICATIONS}
    \hspace{0.5em}
    \rule[0.5ex]{0.395\linewidth}{0.5pt}
\end{center}
\vspace{-0.2em}

\begin{itemize}[leftmargin=0pt, parsep=0pt, itemsep=0.2em]
\item[] \textbf{Bansal, A.}, Mathai, B., Jadhav, N., Balakrishnan, M., Kumar, P. S. (2025). \textbf{\textit{Designing an Algal Façade for Improving Energy Efficiency in a Building in North India}}. Oral presentation at the X International Conference on Sustainable Energy and Environmental Challenges (X SEEC)

\end{itemize}

\vspace{0.3em}

% \section*{Projects}
\noindent
\begin{center}
    \rule[0.5ex]{0.42\linewidth}{0.5pt}
    \hspace{0.5em}
    \textbf{PROJECTS}
    \hspace{0.5em}
    \rule[0.5ex]{0.42\linewidth}{0.5pt}
\end{center}

\begin{itemize}[leftmargin=0pt, parsep=0pt, itemsep=0.8em]

\item[] \textbf{Multi-Agent RL for Linear Salp-Chain Underwater Navigation} \textbar\ \textit{Prof. Sandeep Manjanna} \hfill Jan 2026 - May 2026
{\small
\begin{itemize}[leftmargin=1em, label=\textbullet, itemsep=0.2pt, parsep=0pt, topsep=1pt]
    \item \textbf{Developed a control policy} for a linear chain of underwater robots for coordinated obstacle avoidance navigation. 
    \item Worked with a \textbf{2D underwater navigation} RL environment with salp chain agents from GRASP Lab, UPenn.
    \item Implemented \textbf{multi-agent policies} IQL, MAPPO, and MADDPG for target navigation and obstacle avoidance.
    \item Designed \textbf{dynamic underwater obstacle-placement curriculum}, achieving an \textbf{81\%} collision-free navigation success rate.
\end{itemize}
}

\item[] \textbf{Deep Learning-Enhanced Metabolic Control for PHA Production} \textbar\ \textit{Prof. Monika Sharma} \hfill Aug 2026 - Present
{\small
\begin{itemize}[leftmargin=1em, label=\textbullet, itemsep=0.2pt, parsep=0pt, topsep=1pt]
    \item Building a computational framework to optimize the \textbf{production of bioplastic (PHA)} in Pseudomonas putida.
    \item Developing \textbf{metabolic sub-network reconstruction} for enhancing PHA production pathways.
    \item Formulating a \textbf{dynamic simulations and MCA} based framework to analyze metabolic engineering targets. 
\end{itemize}
}

\item[] \textbf{Inverse PINN Modelling of Flow in a Collapsible Compliant Tube} \textbar\ \textit{Prof. Shashikant Pawar} \hfill Aug 2026 - Present
{\small
\begin{itemize}[leftmargin=1em, label=\textbullet, itemsep=0.2pt, parsep=0pt, topsep=1pt]
    \item Analyzing flow--structure interactions in collapsible tubes to study \textbf{pressure and velocity fields} in venous blood flow.
    \item Constructing \textbf{fluid--structure interaction (FSI) models} and experimental setup for tube deformation and flow measurements.
    \item Implementing a \textbf{physics-informed inverse model (PINN)} for validating experimental measurements.
\end{itemize}
}

\item[] \textbf{Genetic Engineering for Cadmium Phytoremediation} \textbar\ \textit{Prof. Arshdeep Sidhu} \hfill Mar 2026- May 2026
{\small
\begin{itemize}[leftmargin=1em, label=\textbullet, itemsep=0.2pt, parsep=0pt, topsep=1pt]
    \item Designed a genetic engineering strategy for \textit{Populus x canescens} to \textbf{enhance cadmium uptake and stress tolerance.}
    \item Developed a multi-gene construct targeting \textbf{PcNRAMP1, GSH1/GSH2, SOD, and APX} for enhanced phytoremediation.
    \item Proposed framework increases \textbf{cadmium uptake more than 50\%} based on hypothetical estimates. 
\end{itemize}
}

% \item[] \textbf{Microsatellite Instability and WRN Inhibition in Colorectal Cancer} \textbar\ \textit{Prof. Swagata Halder} \hfill Oct 2025 -- Feb 2026
% {\small
% \begin{itemize}[leftmargin=1em, label=\textbullet, itemsep=0.2pt, parsep=0pt, topsep=1pt]
%     \item Conducted a comprehensive literature review on MSI biology and synthetic lethality mechanisms involving WRN helicase.
%     \item Analyzed the therapeutic potential of WRN inhibitors in MSI-high colorectal cancer.
%     \item Working on a review article studying MSI biology, WRN dependency, and its clinical applications.
% \end{itemize}
% }

% \item[] \textbf{Context Aware Cricket Outcomes Prediction} \textbar\ \textit{Prof. Siddharth S.} \hfill Feb 2025 -- May 2025
% {\small
% \begin{itemize}[leftmargin=1em, label=\textbullet, itemsep=0.2pt, parsep=0pt, topsep=1pt]
%     \item Trained a machine learning model to predict an ODI innings' final score, accounting for match context.
%     \item Developed an LSTM-based recurrent neural network aided by convolutional layers to predict innings scorelines.
%     \item Trained the predictor on a custom feature-rich dataset comprising ODI matches since 2006, achieving an RMSE of 13.
% \end{itemize}
% }
\vspace{0.7em}


\item[] \textbf{ZnO-NiO Nanocomposite Synthesis} \textbar\ \textit{Prof. Chaitanya Lekshmi Indira} \hfill Mar 2025 -- May 2025
{\small
\begin{itemize}[leftmargin=1em, label=\textbullet, itemsep=0.2pt, parsep=0pt, topsep=1pt]
    \item Fabricated and analyzed a 1:1 \textbf{ZnO-NiO nanocomposite} using an Ultrasonication-Assisted Mechanical method.
    \item Characterized samples using \textbf{XRD, FTIR, and UV-Vis spectroscopy} and obtained distinctive intensity and absorption peaks.
    \item Achieved up to \textbf{90\% similarity} between sample results and characteristic properties of lab-grade chemicals.
    
\end{itemize}
}

\end{itemize}
    
%  \item[] \textbf{Plaksha EBay} \textbar\ \textit{Prof. Manoj Kannan} \hfill Jan 2024 -- May 2024
% {\small
% \begin{itemize}[leftmargin=1em, label=\textbullet, itemsep=1pt, parsep=1pt, topsep=1pt]
%         \item Designed and implemented a C++ OOP-based platform for student-to-student buying and selling.
%         \item Built an interface to reduce spam,and organize information and space for transactions.
%         \item Developed features for users to streamline campus exchange activities.
%     \end{itemize}
%     }

\end{itemize}

\vspace{0.7em}

% \section*{Skills}
\noindent
\begin{center}
    \rule[0.5ex]{0.44\linewidth}{0.5pt}
    \hspace{0.5em}
    \textbf{SKILLS}
    \hspace{0.5em}
    \rule[0.5ex]{0.44\linewidth}{0.5pt}
\end{center}

\begin{itemize}[leftmargin=1em, label=\textbullet, itemsep=1pt, parsep=1pt, topsep=1pt]
{\small
    \item \textbf{Biosystems:} Microbial culture, molecular cloning, DNA extraction, qPCR amplification, electrochemical analysis
    \item \textbf{Programming Languages:}  C, C++, Python
    \item \textbf{Computational Tools:} BLAST, SnapGene, ClustalW, MEGA; DesignBuilder, COMSOL, OpenLCA}
\end{itemize}

\vspace{0.7em}

%\section*{Relevant Coursework}
\noindent
\begin{center}
    \rule[0.5ex]{0.33\linewidth}{0.5pt}
    \hspace{0.5em}
    \textbf{RELEVANT COURSEWORK}
    \hspace{0.5em}
    \rule[0.5ex]{0.33\linewidth}{0.5pt}
\end{center}

{\small
\noindent
\begin{tabularx}{\textwidth}{@{}X@{\hspace{3em}}l@{\hspace{3em}}X@{}l@{}}
\textbf{BE2002:} Biochemistry and Molecular Biology & Grade: A- &
\textbf{BE3012:} Cell Biology & Grade: A- \\

\textbf{BE3007:} Genetics and Genetic Engineering & Grade: A- &
\textbf{BE2007:} Bioprocess Engineering & Grade: A \\

\textbf{BE3001:} Stochastic Modelling in Biology & Grade: A- &
\textbf{BE2005:} Material Science for Bioengineering & Grade: A- \\

\textbf{BE3014:} Lab-on-a-Chip Devices & Grade: A- &
\textbf{ME3002:} Fluid Mechanics & Grade: A- \\

\textbf{RO3004:} Reinforcement Learning & Grade: A- &
\textbf{AI3011:} Machine Learning & Grade: A
\end{tabularx}
}

\vspace{0.2em}

% \section*{Leadership \& Roles}
\noindent
\begin{center}
    \rule[0.5ex]{0.36\linewidth}{0.5pt}
    \hspace{1em}
    \textbf{LEADERSHIP ROLES}
    \hspace{1em}
    \rule[0.5ex]{0.36\linewidth}{0.5pt}
\end{center}

\begin{itemize}[leftmargin=0pt, parsep=0pt, itemsep=0.8em]

\item[] \textbf{Cohort Representative} \textbar\ {Biological Systems Engineering, Plaksha University}  \hfill Jan 2025 -- Present
{\small
\begin{itemize}[leftmargin=1em, label=\textbullet, itemsep=0.2pt, parsep=0pt, topsep=1pt]
    \item Representing BSE cohort, facilitating communication between 20+ faculty members and the student body.
    \item Bridging student‐faculty communication on course policies resolving student academic challenges through continuous feedback.
\end{itemize}
}

\item[] \textbf{President} \textbar\ {Nebula, Plaksha Well-Being Club} \hfill Feb 2025 -- Present
{\small
\begin{itemize}[leftmargin=1em, label=\textbullet, itemsep=0.2pt, parsep=0pt, topsep=1pt]
    \item Leading a 30 member team to design and implement student-centric mental health and stress-relief programs.
    \item Conducted events and competitions to contribute to a more supportive and resilient campus culture.
\end{itemize}
}

\item[] \textbf{Core Member, Organizing Committee} \textbar\ {Dotslash Hackathon, Plaksha University} \hfill Nov 2024
{\small
\begin{itemize}[leftmargin=1em, label=\textbullet, itemsep=0.2pt, parsep=0pt, topsep=1pt]
    \item Collaborated on logistics, venue layout, and challenge-track design for a 120+ participant inter-college hackathon.
\end{itemize}
}

\item[] \textbf{House Captain} \textbar\ {MGD Girls' School} \hfill Jul 2022 -- Apr 2023
{\small
\begin{itemize}[leftmargin=1em, label=\textbullet, itemsep=0.2pt, parsep=0pt, topsep=1pt]
    \item Led 500+ students in cultural and sports events, securing the award for best house in academics and co-curricular activities.
    \item Coordinated major annual events for 1000+ students as part of a 37-member student council.
\end{itemize}
}

\end{itemize}

    % \vspace{0.2em}
% \begin{itemize}[leftmargin=*]
%     \item textbf{Class Representative, BSE Cohort}
    
%     \item Organising committee for inter-college hackathon (Dotslash)
%     \item Finance lead for Rangrez (Plaksha Art Club)
%     \item Co-president for Nebula (Plaksha Well-being Club)
%     \item House captain in school – managed cultural and sports events


\vspace{0.7em}

% \section*{Achievements}
\noindent
\begin{center}
    \rule[0.5ex]{0.38\linewidth}{0.5pt}
    \hspace{1em}
    \textbf{ACHIEVEMENTS}
    \hspace{1em}
    \rule[0.5ex]{0.38\linewidth}{0.5pt}
\end{center}

\begin{itemize}[leftmargin=1em, label=\textbullet, itemsep =0pt, parasep=0pt]
        \item \textbf{Runners-up, Hack Plaksha} (Fintech Track, Feb 2024) - Developed Dhanshakti, a financial management app to empower rural self-help groups by simplifying savings and loan tracking.
        \item \textbf{Gold Medal, District Karate Tournament} - Secured first place in Individual Shihai in Under 18 category.
        \item \textbf {Merit Scholarship, Plaksha University}: Recognized for academic excellence with a GPA above 8 and extracurricular.
        \item Completed Foundation certification for \textbf{BS in Data Science} from IIT Madras.

    \end{itemize}
\end{document}